\documentclass[10pt]{article}
\usepackage[T1]{fontenc}
\usepackage[utf8]{inputenc}
\usepackage[french]{babel}
\usepackage{amsmath,amssymb,mathrsfs,mathpazo}
\usepackage{geometry}
\geometry{a4paper,margin=2cm}
\usepackage{graphicx,array,tabularx,tabularray,colortbl,multirow,multicol,booktabs,enumitem,pifont,ulem,fancybox,diagbox,cancel}
\usepackage{pstricks,pstricks-add,pst-plot,pst-tree,pst-node,pst-text,pst-func,pst-math,pst-bspline,pst-3dplot,pst-eucl,pst-solides3d}
\let\pscalculate\undefined
\usepackage{pst-bezier}
\usepackage{tikz,pgfplots,tkz-tab}
\usetikzlibrary{arrows,arrows.meta,patterns,calc,positioning,trees,decorations.markings,shapes.geometric,shapes.misc,shapes.arrows,intersections,decorations.pathmorphing,decorations.pathreplacing,backgrounds,shadows,babel}
\pgfplotsset{compat=1.18}
\usepgfplotslibrary{fillbetween}
\usepackage[np]{numprint}
\usepackage{listings,makecell,pas-tableur,twoopt}
\let\barre\undefined
\usepackage[pstricks]{tablvar}
\definecolor{darkgreen}{RGB}{0,100,0}
\providecommand{\red}{\color{red}}
\providecommand{\blue}{\color{blue}}
\providecommand{\green}{\color{green}}
\providecommand{\black}{\color{black}}
\providecommand{\white}{\color{white}}
\providecommand{\gray}{\color{gray}}
\providecommand{\cyan}{\color{cyan}}
\providecommand{\magenta}{\color{magenta}}
\providecommand{\yellow}{\color{yellow}}
\providecommand{\violet}{\color{violet}}
\providecommand{\orange}{\color{orange}}
\providecommand{\euro}{\texteuro}
\providecommand{\EUR}{\texteuro}
\providecommand{\ieme}{\textsuperscript{e}}
\providecommand{\degres}{\ensuremath{^\circ}}
\providecommand{\degree}{\ensuremath{^\circ}}
\providecommand{\vv}[1]{\overrightarrow{#1}}
\providecommand{\Coord}[1]{#1}
\providecommand{\decosix}{}
\providecommand{\starredbullet}{\ensuremath{\bullet}}
\pagestyle{empty}
\setlength{\parindent}{0pt}
\usepackage{adjustbox}\DeclareUnicodeCharacter{2236}{:}\DeclareUnicodeCharacter{2190}{\ensuremath{\leftarrow}}
\providecommand{\R}{\mathbb{R}}
\providecommand{\N}{\mathbb{N}}
\providecommand{\Z}{\mathbb{Z}}
\providecommand{\Q}{\mathbb{Q}}
\providecommand{\D}{\mathbb{D}}
\providecommand{\C}{\mathbb{C}}
\providecommand{\e}{{\mathrm{e}}}
\providecommand{\vect}[1]{\overrightarrow{#1}}
\providecommand{\vectt}[1]{\overrightarrow{\mathrm{#1}}}
\providecommand{\vvt}[1]{\overrightarrow{\mathrm{#1}}}
\providecommand{\np}[2][]{{#2~#1}}
\providecommand{\vv}[1]{{\overrightarrow{#1}}}
\renewcommand{\d}{{\mathrm{d}}}
\providecommand{\degres}{{^\circ}}
\providecommand{\degre}{{^\circ}}
\providecommand{\blue}{\color{blue}}
\providecommand{\red}{\color{red}}
\providecommand{\enskip}{\quad}
\providecommand{\llbracket}{\lbrack\!\lbrack}
\providecommand{\rrbracket}{\rbrack\!\rbrack}
\providecommand{\euro}{{\text{€}}}
\providecommand{\no}{{\text{n°}}}
\providecommand{\interval}[2]{{[#1;#2]}}
\providecommand{\intervFO}[2]{{[#1;#2[}}
\providecommand{\dotfill}{\ldots}

\providecommand{\ds}{\displaystyle}
\providecommand{\textding}[1]{#1}
\providecommand{\Oij}{\(\mathrm{O};\vec{\imath},\vec{\jmath}\)}
\providecommand{\Oijk}{\(\mathrm{O};\vec{\imath},\vec{\jmath},\vec{k}\)}
\providecommand{\Ouv}{\(\mathrm{O};\vec{u},\vec{v}\)}
\providecommand{\intervFF}[2]{\ensuremath{[#1;#2]}}
\renewcommand{\euro}{\ifmmode\text{\texteuro}\else\texteuro\fi}
\DeclareUnicodeCharacter{00B0}{\ifmmode^{\circ}\else\textdegree\fi}
\DeclareUnicodeCharacter{20AC}{\ifmmode\text{\texteuro}\else\texteuro\fi}
% Énoncés : Éducation nationale. Transcription et corrigés : APMEP, Denis Vergès et contributeurs. Découpage et classement : CoopMaths, Sébastien Lozano, Éric Elter et contributeurs. Adaptation de la présentation : Maths Explorer.
% CC BY-SA : https://coopmaths.fr/www/about/
% https://coopmaths.fr/alea/static/e3c/2020/tex/e3c_2020_00_sujet14_1_cor.tex
\begin{document}
\noindent\textit{APMEP / CoopMaths. CC BY-SA. Présentation adaptée par Maths Explorer.}\par\medskip
\begin{enumerate}
		\item 	L'inéquation \(x^2 + x + 2 > 0\) :
		
		\[
		x^2 + x + 2 > 0 \iff \left( x + \dfrac{1}{2} \right)^2 - \dfrac{1}{4} + 2 > 0 \iff \left( x + \dfrac{1}{2} \right)^2 + \dfrac{7}{4} > 0.
		\]
		
		Comme \(\dfrac{7}{4} > 0\), ceci est vrai quel que soit le réel \(x\) : \(S = \mathbb{R}\).
		
		\item 
		
		Soient \(\overrightarrow{u}\) et \(\overrightarrow{v}\) deux vecteurs tels que \(\|\overrightarrow{u} \| = 3\), \(\|\overrightarrow{v} \| = 2\) et \(\overrightarrow{u} \cdot \overrightarrow{v} = -1\),\\
		alors \(\|\overrightarrow{u} + \overrightarrow{v} \|^2\) est égal à :
		
		\[
		\|\overrightarrow{u} + \overrightarrow{v} \|^2 = (\overrightarrow{u} + \overrightarrow{v}) \cdot (\overrightarrow{u} + \overrightarrow{v}) = \overrightarrow{u} \cdot \overrightarrow{u} + \overrightarrow{v} \cdot \overrightarrow{v} + 2 \overrightarrow{u} \cdot \overrightarrow{v} = 3^2 + 2^2 + 2 \times (-1) = 9 + 4 - 2 = 11.
		\]
		
		\item 
		
		Soient \(A\) et \(B\) deux évènements d'un univers tels que \(P_A(B) = 0,2\) et \(P(A) = 0,5\).\\ Alors la probabilité \(P(A \cap B)\) est égale à :
		
		\[
		P_A(B) = \dfrac{P(A \cap B)}{P(A)} \iff 0,2 = \dfrac{P(A \cap B)}{0,5} \Rightarrow P(A \cap B) = 0,2 \times 0,5 = 0,1.
		\]
		
		\item 
		
		Soit \((u_n)\) une suite arithmétique de terme initial \(u_0 = 2\) et de raison $3$.\\ 
		La somme \(S\) définie par \(S = u_0 + u_1 + \ldots + u_{12}\) est égale à :
		
		\[
		S = 2 + 5 + 8 + 11 + \ldots + 35 + 38
		\]
		
		ou
		
		\[
		S = 38 + 35 + 32 + \ldots + 8 + 5 + 2.
		\]
		
		En faisant la somme, on obtient :
		
		\[
		2S = 40 + 40 + 40 + \ldots + 40 = 13 \times 40 = 520 \Rightarrow S = 260.
		\]
		
		\item 
		
		On a \(f'(x) = 3(2x - 5)^2 \times (2x - 5)' = 3 \times 2(2x - 5)^2 = 6(2x - 5)^2\).
	\end{enumerate}

\end{document}
